\section{From slider SER to variational empirical Bayes}
\label{sec:slider-veb}

The one-SNP calculation extends to a single-effect regression (SER)
with uncertainty over which SNP has an effect. Fitting this SER model
by empirical Bayes (EB) to the expected residuals of the other effects
then gives a block coordinate-ascent algorithm for a variational
empirical-Bayes (VEB) objective. The argument follows the general
additive-effects construction of Wang et al.~\cite{wang2018susie},
Appendix B, especially Proposition 2. Here ``additive effects'' means
that the component contributions are summed; each component may have
a nonadditive genotype coding.

The distinction between SNP-specific and shared sliders is essential.
Within one SER component we fit one slider for each candidate SNP.
In the sum of $L$ components these parameters are
$\Delta=(\delta_{\ell j})$, with a separate slider for every
component--SNP pair. They are model parameters estimated by EB.
The latent variables are the selected SNP identities and their effect
sizes.

\subsection{Working coordinates and a slider SER model}

We condition on the observed genotype matrix. To retain the
intercept-integrated convention used above, let
$\mathbf U\in\mathbb R^{n\times N}$ have orthonormal columns spanning
the orthogonal complement of $\mathbf 1$, where $N=n-1$. Thus
$\mathbf U^\top\mathbf U=I_N$ and
$\mathbf U\mathbf U^\top=I_n-\mathbf 1\mathbf 1^\top/n$.
Define
\begin{equation}
\widetilde{\mathbf y}=\mathbf U^\top\mathbf y,\qquad
\mathbf z_j(\delta)
 =\frac{\mathbf U^\top(\mathbf x_j+\delta\mathbf h_j)}{s_j},
\qquad h_{ij}=\ind_{\{x_{ij}=1\}}.
\label{eq:veb-working-design}
\end{equation}
No explicit construction of $\mathbf U$ is needed: inner products in
these coordinates equal the corresponding centered inner products.
Without an intercept take $\mathbf U=I_n$ and $N=n$.
The scales $s_j>0$ are fixed functions of the genotypes, independent
of $\delta$. Taking $s_j=1$ gives the raw-scale model above; taking
$s_j$ to be the additive genotype standard deviation matches the
scaling convention in \texttt{susieSlide}. Constant predictors are
excluded when this latter convention is used. The Gaussian effect
prior is defined on the chosen scale.

For a response $\mathbf r\in\mathbb R^N$, the slider SER model is
\begin{align}
J&\sim\operatorname{Categorical}(\pi_1,\ldots,\pi_p),
 &\beta&\sim\Normal(0,V),\qquad J\perp\beta,\nonumber\\
\mathbf r\mid J=j,\beta,\boldsymbol\delta
 &\sim\Normal\{\beta\mathbf z_j(\delta_j),\sigma^2 I_N\},
 &\delta_j&\in\mathcal D_j.
\label{eq:veb-ser-model}
\end{align}
Here $V>0$, $\sigma^2>0$, and the prior weights are fixed, with
$\pi_j>0$ and $\sum_j\pi_j=1$. Ordinarily
$\mathcal D_j=[-1,1]$; a genotype-count rule can instead impose
$\mathcal D_j=\{0\}$. These constraints depend only on the observed
genotypes. There is exactly one selected SNP in this component,
although a slider is specified for every candidate SNP.

For any fixed $\boldsymbol\delta$, put
\begin{align}
S_j(\delta)&=\|\mathbf z_j(\delta)\|^2,
 &T_j(\delta)&=\mathbf z_j(\delta)^\top\mathbf r,\nonumber\\
v_j(\delta)&=\left\{V^{-1}+S_j(\delta)/\sigma^2\right\}^{-1},
 &m_j(\delta)&=v_j(\delta)T_j(\delta)/\sigma^2.
\label{eq:veb-ser-moments}
\end{align}
Integration over $\beta$ gives
\begin{equation}
\BF_j(\delta;V,\sigma^2)
 =\left\{1+\frac{VS_j(\delta)}{\sigma^2}\right\}^{-1/2}
 \exp\left\{\frac{VT_j(\delta)^2}
 {2\sigma^2\{\sigma^2+VS_j(\delta)\}}\right\}.
\label{eq:veb-ser-bf}
\end{equation}
Thus, with $p_0(\mathbf r;\sigma^2)$ the $\Normal(0,\sigma^2 I_N)$
density, the SER marginal likelihood and exact conditional posterior
are
\begin{align}
M(\mathbf r;\boldsymbol\delta,V,\sigma^2)
 &=p_0(\mathbf r;\sigma^2)
   \sum_{j=1}^p\pi_j\BF_j(\delta_j;V,\sigma^2),
\label{eq:veb-ser-evidence}\\
a_j:=\Pr(J=j\mid\mathbf r,\boldsymbol\delta)
 &=\frac{\pi_j\BF_j(\delta_j;V,\sigma^2)}
 {\sum_{k=1}^p\pi_k\BF_k(\delta_k;V,\sigma^2)},
\label{eq:veb-ser-weights}\\
\beta\mid J=j,\mathbf r,\boldsymbol\delta
 &\sim\Normal\{m_j(\delta_j),v_j(\delta_j)\}.
\label{eq:veb-ser-posterior}
\end{align}
The variances are suppressed in the posterior conditioning notation.

\subsection{Why independent slider maximizations solve the SER EB problem}

\paragraph{Proposition 1 (SNP-specific slider EB update).}
For fixed $V$, $\sigma^2$, and $\boldsymbol\pi$, choose
\begin{equation}
\widehat\delta_j\in\argmax_{\delta\in\mathcal D_j}
       \BF_j(\delta;V,\sigma^2),\qquad j=1,\ldots,p.
\label{eq:veb-separate-maxima}
\end{equation}
Then $\widehat{\boldsymbol\delta}$ maximizes the SER marginal
likelihood in equation~\eqref{eq:veb-ser-evidence} over
$\mathcal D_1\times\cdots\times\mathcal D_p$.

\paragraph{Proof.}
Each summand in equation~\eqref{eq:veb-ser-evidence} depends on only
one slider. For every admissible $\boldsymbol\delta$,
\begin{equation}
\sum_j\pi_j\BF_j(\delta_j;V,\sigma^2)
 \leq\sum_j\pi_j\BF_j(\widehat\delta_j;V,\sigma^2).
\label{eq:veb-separability-proof}
\end{equation}
The null density is positive and independent of
$\boldsymbol\delta$, so the same inequality holds for $M$.
Continuity on the compact allowed intervals ensures that maximizers
exist. This proves the result.\hfill$\square$

The cubic maximization already derived in equation~\eqref{eq:cubic}
therefore solves each required one-dimensional problem. No posterior
SNP weight is needed in that maximization: the evidence is a sum of
independently adjustable nonnegative terms. The posterior weights
are recomputed from equation~\eqref{eq:veb-ser-weights} afterwards.
If a \emph{single} slider were shared across candidate SNPs, the
correct objective would instead be
$\sum_j\pi_j\BF_j(\delta;V,\sigma^2)$; separate candidate-wise
maximizations would not solve that different model.

\subsection{The sum-of-single-effects model and its ELBO}

Let the $L$ pairs $(J_\ell,\beta_\ell)$ be independent a priori,
with $J_\ell\sim\operatorname{Categorical}(\boldsymbol\pi)$ and
$\beta_\ell\sim\Normal(0,V_\ell)$. Define
\begin{equation}
\boldsymbol\eta_\ell
 =\beta_\ell\mathbf z_{J_\ell}(\delta_{\ell J_\ell}),\qquad
\widetilde{\mathbf y}=\sum_{\ell=1}^L\boldsymbol\eta_\ell
      +\boldsymbol\epsilon,\qquad
\boldsymbol\epsilon\sim\Normal(0,\sigma^2 I_N).
\label{eq:veb-sum-model}
\end{equation}
For every fixed $\Delta$, this specifies a proper probability model
on the working coordinates. Equivalently, the prior on
$\boldsymbol\eta_\ell$ is a mixture over $p$ Gaussian single-effect
directions, with directions determined by the sliders. This is a
special case of the general additive-effects model in Appendix B.2
of Wang et al.~\cite{wang2018susie}.

We use the variational family
\begin{equation}
q(J_1,\beta_1,\ldots,J_L,\beta_L)
  =\prod_{\ell=1}^Lq_\ell(J_\ell,\beta_\ell).
\label{eq:veb-family}
\end{equation}
Each factor is unrestricted on the pair $(J_\ell,\beta_\ell)$;
it does not factorize over candidate SNPs. Write $p_\ell$ for the
categorical--Gaussian prior on this pair and
$\theta=(\Delta,V_1,\ldots,V_L,\sigma^2)$. The objective is
\begin{align}
\mathcal F(q,\theta)
 &=E_q\log p(\widetilde{\mathbf y}\mid J,\beta,\Delta,\sigma^2)
   -\sum_{\ell=1}^L\operatorname{KL}(q_\ell\|p_\ell)
\label{eq:veb-elbo}\\
 &=\log p(\widetilde{\mathbf y}\mid\theta)
   -\operatorname{KL}\{q\|p(J,\beta\mid
                           \widetilde{\mathbf y},\theta)\}.
\label{eq:veb-bound-identity}
\end{align}
Consequently $\mathcal F$ is a lower bound on the log marginal
likelihood for every fixed $\theta$. Maximizing it over the
factorized family and over the chosen model parameters is a VEB
optimization problem. No prior on $\Delta$ is required for this
interpretation: $\Delta$ is optimized as a parameter, whereas
$q$ approximates the posterior of the latent variables conditional
on those parameter estimates.

\subsection{A slider SER update is an exact VEB block update}

\paragraph{Proposition 2 (residual SER update).}
Hold $\sigma^2$, $q_k$, $V_k$, and $\boldsymbol\delta_k$ fixed
for all $k\neq\ell$, and define the expected residual
\begin{equation}
\mathbf r_\ell=\widetilde{\mathbf y}
       -\sum_{k\neq\ell}\overline{\boldsymbol\eta}_k,
\qquad
\overline{\boldsymbol\eta}_k=E_{q_k}\boldsymbol\eta_k.
\label{eq:veb-expected-residual}
\end{equation}
For fixed $V_\ell$, jointly maximizing $\mathcal F$ over
$(q_\ell,\boldsymbol\delta_\ell)$ is achieved by fitting the
slider SER model to $\mathbf r_\ell$: first apply
equation~\eqref{eq:veb-separate-maxima}, then set $q_\ell$ to the
exact posterior in equations~\eqref{eq:veb-ser-weights}--\eqref{eq:veb-ser-posterior}.

\paragraph{Proof.}
Independence across components under $q$ gives
\begin{align}
E_q\left\|\widetilde{\mathbf y}
          -\sum_k\boldsymbol\eta_k\right\|^2
 &=E_{q_\ell}\|\mathbf r_\ell-\boldsymbol\eta_\ell\|^2
   +C_{-\ell},\nonumber\\
C_{-\ell}
 &=\sum_{k\neq\ell}
   E_{q_k}\|\boldsymbol\eta_k
                 -\overline{\boldsymbol\eta}_k\|^2.
\label{eq:veb-residual-decomposition}
\end{align}
The second term does not depend on the parameters of component
$\ell$ or on $q_\ell$. Thus
\begin{align}
\mathcal F(q,\theta)
 &=C+\mathcal F_{\mathrm{SER}}
       (q_\ell,\boldsymbol\delta_\ell,V_\ell;\mathbf r_\ell)
\label{eq:veb-block-identity}\\
 &=C+\log M(\mathbf r_\ell;
           \boldsymbol\delta_\ell,V_\ell,\sigma^2)
   -\operatorname{KL}(q_\ell\|p_{\mathrm{SER},\ell}),
\label{eq:veb-block-kl}
\end{align}
where $C$ is independent of $(q_\ell,\boldsymbol\delta_\ell,V_\ell)$,
and $p_{\mathrm{SER},\ell}$ is the exact SER posterior at the
displayed parameters. For any fixed slider vector, the KL term
has minimum zero, attained at that posterior. It follows that
\begin{equation}
\max_{q_\ell,\boldsymbol\delta_\ell}\mathcal F(q,\theta)
 =C+\max_{\boldsymbol\delta_\ell}
       \log M(\mathbf r_\ell;
                \boldsymbol\delta_\ell,V_\ell,\sigma^2).
\label{eq:veb-profiled-block}
\end{equation}
Proposition 1 solves the remaining maximization. This proves the
claim.\hfill$\square$

This is a \emph{joint} update of the slider vector and its
variational factor. It need not increase the ELBO if one changes
the sliders while keeping the old posterior moments. The posterior
weights, means, variances, and fitted component must all be refreshed
at the new slider values. Also, the residual construction does not
assume that the other components are known: their posterior
uncertainty is the term $C_{-\ell}$, which is constant only for this
component update at fixed noise variance.

\subsection{Algorithm, variance updates, and interpretation}

Cycle over components. For each $\ell$, compute
$\mathbf r_\ell$, maximize every candidate's log-BF over
$\mathcal D_j$, compute the exact SER posterior at those sliders,
and refresh
\begin{equation}
\overline{\boldsymbol\eta}_\ell
 =\sum_j a_{\ell j}m_{\ell j}
              \mathbf z_j(\delta_{\ell j}).
\label{eq:veb-component-mean}
\end{equation}
Here $a_{\ell j}$, $m_{\ell j}$, and $v_{\ell j}$ denote the
current posterior probability, conditional mean, and conditional
variance. With exact block maximization each update is
nondecreasing in the same ELBO. This establishes a valid VEB
coordinate-ascent algorithm; it does not establish a global optimum.
The local SER problem is solved exactly conditional on its EB
parameters, while the approximation in the full model is the
factorization across components.

If $V_\ell$ is also learned by marginal likelihood, profile the
sliders at each trial value of $V$ and optimize
\begin{equation}
\widehat V_\ell\in\argmax_V
 \log\left\{\sum_j\pi_j
  \BF_j\bigl(\widehat\delta_{\ell j}(V);V,\sigma^2\bigr)\right\}.
\label{eq:veb-profile-prior-variance}
\end{equation}
The variance is shared across the candidate SNPs in that component,
so this outer maximization does not separate over SNPs. A numerical
outer search should retain the previous value unless it improves
the objective. Alternatively, an EM coordinate update at fixed
$q$ is $V_\ell\leftarrow\sum_j a_{\ell j}(v_{\ell j}+m_{\ell j}^2)$,
followed by subsequent posterior updates. The derivation above
assumes positive variances; $V_\ell=0$ can be handled separately
as an inactive, zero-effect component.

The expected residual sum of squares is
\begin{align}
\operatorname{ERSS}
 &=\left\|\widetilde{\mathbf y}
        -\sum_\ell\overline{\boldsymbol\eta}_\ell\right\|^2
   +\sum_\ell\left\{
      \sum_j a_{\ell j}(v_{\ell j}+m_{\ell j}^2)
                    S_j(\delta_{\ell j})
      -\|\overline{\boldsymbol\eta}_\ell\|^2\right\}.
\label{eq:veb-erss}
\end{align}
At fixed $q$, $\Delta$, and prior variances, maximizing the ELBO
over the noise variance gives $\sigma^2\leftarrow\operatorname{ERSS}/N$.
The dimension is $N=n-1$ for the intercept-integrated likelihood used
here. An implementation that instead estimates the intercept as a
fixed parameter in the full Gaussian likelihood uses $n$ in the
normalizing term and this denominator, as does the current
\texttt{susieSlide} implementation. The two conventions give the
same centered SER updates at a fixed $\sigma^2$, but their
noise-variance updates should not be conflated.

For completeness, a directly evaluable expression for the ELBO is
\begin{align}
\mathcal F
 &=-\frac N2\log(2\pi\sigma^2)
   -\frac{\operatorname{ERSS}}{2\sigma^2}
   -\sum_\ell K_\ell,\nonumber\\
K_\ell
 &=\sum_j a_{\ell j}\log\frac{a_{\ell j}}{\pi_j}
  +\frac12\sum_j a_{\ell j}
   \left\{\frac{v_{\ell j}+m_{\ell j}^2}{V_\ell}
                -1+\log\frac{V_\ell}{v_{\ell j}}\right\}.
\label{eq:veb-computable-elbo}
\end{align}
This expression includes both uncertainty about SNP identity and
uncertainty about the effect size.

For active components, the variational probability that SNP $j$
is selected by at least one component is
\begin{equation}
\operatorname{PIP}_j
  =1-\prod_{\ell=1}^L(1-a_{\ell j}).
\label{eq:veb-pip}
\end{equation}
A component credible set collects SNPs whose $a_{\ell j}$ sum to
at least the requested probability. Both summaries condition on
the fitted sliders and variances. The VEB derivation establishes
the model and objective; frequentist coverage and PIP calibration
remain empirical properties to assess. In particular,
$\BF_j(\widehat\delta_{\ell j})$ is conditional fitted-slider
evidence, not evidence averaged over a prior for $\delta$.

The separate-component assumption matters. A slider shared across
all components for a given SNP would change several component
contributions simultaneously and require a different parameter
update; the candidate-wise update above would not by itself prove
ascent for that model. In the present model, multiple components
may select the same SNP with different sliders. The monotonicity
constraint applies to each single effect; their sum need not be
monotone if their effect directions differ. Setting all sliders to
zero recovers additive SuSiE on the same working design, with the
same priors and variance conventions.
